% ============================================================
% SECTION 8 — REFUND STATE MACHINE
% ============================================================

\section{Refund State Machine}

\begin{tcolorbox}[summarybox={\iconFlow[1.5] Reversing the Transaction}]
\color{textdark}
While the \code{RecoveryAttempt} state machine governs charging money, the \textbf{Refund State Machine} governs returning it. A refund is arguably a higher-risk operation: double-charging a customer results in a support ticket; double-refunding a customer results in direct, unrecoverable capital loss for the merchant.

The Refund State Machine is strictly decoupled from the recovery state, living directly on the \code{Transaction} model via the \code{refund\_status} field.
\end{tcolorbox}

\vspace{0.8cm}

% --- THE REFUND STATE GRAPH ---
\subsection{The Refund State Graph}

Unlike recovery attempts, which generate multiple child records, a refund is a singular operation against the parent \code{Transaction}. The state graph is simpler but enforces strict entry conditions.

\vspace{0.5cm}

\begin{center}
\begin{tikzpicture}[
    node distance=2cm and 1.5cm,
    state/.style={rectangle, draw=primary, thick, fill=bglight, minimum width=3cm, minimum height=0.8cm, rounded corners=4pt, font=\small\bfseries, drop shadow=black!5, align=center},
    terminal/.style={rectangle, draw=textdark, thick, fill=textdark!10, minimum width=3cm, minimum height=0.8cm, rounded corners=4pt, font=\small\bfseries, drop shadow=black!5, align=center},
    ambiguous/.style={rectangle, draw=warning, thick, fill=warning!10, minimum width=3cm, minimum height=0.8cm, rounded corners=4pt, font=\small\bfseries, drop shadow=black!5, align=center},
    success/.style={rectangle, draw=secondary, thick, fill=secondary!10, minimum width=3cm, minimum height=0.8cm, rounded corners=4pt, font=\small\bfseries, drop shadow=black!5, align=center},
    fail/.style={rectangle, draw=danger, thick, fill=danger!10, minimum width=3cm, minimum height=0.8cm, rounded corners=4pt, font=\small\bfseries, drop shadow=black!5, align=center},
    arrow/.style={-{Stealth[scale=1.1]}, thick, draw=textdark},
    label/.style={font=\scriptsize\itshape, text=textmuted}
]

    % States
    \node[state] (none) {NONE};
    \node[state, right=of none] (req) {REFUND\_REQUESTED};
    \node[state, below=1cm of req] (proc) {REFUND\_PROCESSING};
    
    \node[ambiguous, below left=1.2cm and 0cm of proc] (unk) {REFUND\_UNKNOWN};
    \node[success, below=1.2cm of proc] (succ) {REFUNDED};
    \node[fail, below right=1.2cm and 0cm of proc] (fail) {REFUND\_FAILED};

    % Transitions
    \draw[arrow] (none) -- node[above, font=\scriptsize, align=center] {API\\Request} (req);
    \draw[arrow] (req) -- (proc);
    
    \draw[arrow] (proc) -- (succ);
    \draw[arrow] (proc) -- (fail);
    \draw[arrow] (proc) -| (unk);
    
    \draw[arrow] (unk.west) -| node[above left, font=\scriptsize, align=center] {Webhook /\\Reconciliation} (succ.west);
    \draw[arrow] (unk.east) -| node[above right, font=\scriptsize, align=center] {Webhook /\\Reconciliation} (fail.east);

\end{tikzpicture}
\end{center}

\vspace{0.5cm}

% --- ENTRY VALIDATION ---
\subsection{Entry Validation: The "Recovered Only" Rule}

Before a transaction can even enter the \refundstate{REFUND\_REQUESTED} state, \filepath{app/services/refund\_service.py} enforces a strict logical invariant:

\begin{tcolorbox}[warningbox={You cannot refund a failure}]
The system explicitly rejects any refund request unless the underlying \code{Transaction.status} is \code{"success"} OR the \code{Transaction.recovery\_status} is \code{"SUCCEEDED"}. Attempting to refund a failed payment immediately returns an HTTP 400 Bad Request.
\end{tcolorbox}

\vspace{0.8cm}

% --- THE P0 VULNERABILITY (BATCH 4) ---
\subsection{The Batch 4 Vulnerability: The Concurrent Refund Race}

During the Final Adversarial Audit, a \textbf{P0 severity vulnerability} was discovered in the refund initiation logic. The original code looked like this:

\begin{tcolorbox}[codebox={Original Vulnerable Logic (\pzero)}]
\begin{lstlisting}[language=Python]
def initiate_refund(db: Session, transaction_id: str):
    txn = db.query(Transaction).filter_by(id=transaction_id).first()
    
    # VULNERABILITY: Read occurs without a lock
    if txn.refund_status in ["REFUND_REQUESTED", "REFUND_PROCESSING", "REFUNDED"]:
        raise ValueError("Refund already in progress")
        
    # RACE CONDITION: Two threads can reach here simultaneously
    txn.refund_status = "REFUND_REQUESTED"
    db.commit()
    
    # Both threads proceed to call the gateway...
\end{lstlisting}
\end{tcolorbox}

If an attacker (or a buggy merchant UI) fired two simultaneous refund API calls with \textit{different} Idempotency-Keys, the API Idempotency layer (Pillar 1) would allow both through. Both threads would read \code{refund\_status == "NONE"} at the exact same millisecond. Both would pass the check, and both would issue a refund to the gateway, double-refunding the transaction.

\vspace{0.8cm}

% --- THE BATCH 4.6 REMEDIATION ---
\subsection{The Batch 4.6 Remediation: Optimistic State Enforcement}

To fix this, the engineering team could have used a pessimistic lock (\code{with\_for\_update()}). However, because SQLite (used for testing and local dev) silently ignores row locks, the team opted for an \textbf{Optimistic Concurrency Control (OCC)} approach that works robustly across all SQL databases.

\begin{tcolorbox}[codebox={Remediated OCC Logic (Batch 4.6)}]
\begin{lstlisting}[language=Python]
def initiate_refund(db: Session, transaction_id: str):
    # Optimistic Update: Force the DB engine to evaluate the prior state
    updated_rows = db.query(Transaction).filter(
        Transaction.id == transaction_id,
        (Transaction.refund_status == "NONE") | (Transaction.refund_status.is_(None))
    ).update({"refund_status": "REFUND_REQUESTED"})
    
    if updated_rows == 0:
        # The row was modified by another thread between read and write!
        db.rollback()
        raise ValueError("Concurrent refund detected or already refunded.")
        
    db.commit()
\end{lstlisting}
\end{tcolorbox}

By moving the state check into the \code{WHERE} clause of the \code{UPDATE} statement, the system relies on the database engine's internal atomic write lock. Whichever thread arrives first will update the row; the second thread will find 0 rows matching \code{refund\_status == "NONE"} and safely abort.

\vspace{0.8cm}

% --- WEBHOOK INTEGRATION ---
\subsection{Webhook Validation for Refunds}

Because refunds take time (days) to clear banking networks, they heavily rely on asynchronous webhooks from the gateway (\code{refund.processed}, \code{refund.failed}).

Prior to Batch 4.6, the webhook handler blindly accepted the gateway's status. This created an attack vector where a forged or buggy webhook could transition a \code{NONE} transaction directly to \code{REFUNDED}. 

The webhook logic (\filepath{app/api/webhooks.py}) now strictly validates intent:
\begin{enumerate}
    \item When a \code{refund.processed} webhook arrives...
    \item The system checks if \code{txn.refund\_status} is \code{REFUND\_REQUESTED} or \code{REFUND\_PROCESSING}.
    \item If yes, it transitions to \code{REFUNDED}.
    \item If no, it \textbf{rejects} the webhook. The system will never acknowledge a refund it did not explicitly initiate.
\end{enumerate}
